\paragraph{What this paper shows.} A single-shot supervised fine-tune of a 7B
VLM, using bbox supervision distilled from a 72B teacher, can match (or beat)
off-the-shelf 7B detection on the Dish-Correct semantic metric without any
exotic loss function or model surgery. Most of the lift comes from the
\emph{count-exact} term --- the off-the-shelf 7B is biased toward
under-detection, and supervised exposure to multi-item trays is sufficient
to recalibrate its prior. Mean-pooled prototype heads in Stage~3 are too
coarse for our long-tailed taxonomy; max-similarity retrieval over the full
train-split reference bank is a one-line change that better handles
multimodal classes (rice variants, cake variants).

\paragraph{What this paper does not show.} We deliberately scoped the
empirical comparison to two arms (off-the-shelf vs distilled 7B) under
identical Stage 2 + 3, to isolate the effect of distillation. Several
ablations are deferred:

\begin{itemize}\itemsep1pt
  \item \textbf{Cardinality-aware retrain.} The structural count loss
    (\S\ref{sec:cardinality}) is in epoch~$1$ of $10$ at submission time;
    its end-of-training Dish-Correct number is unavailable. The
    in-progress training curves (Figure~\ref{fig:training-curves}) suggest
    the val multiset-match rate is recovering from the regression seen in
    \texttt{bboxonly-epoch6}.
  \item \textbf{DINOv3.} The Stage 3 encoder remains DINOv2-large.
    DINOv3 (released after our experiments started) offers stronger
    embeddings on natural-image benchmarks; swapping it in is a one-line
    config change and would re-embed the $\sim$3{,}900-reference bank in
    $\sim$2~min. We expect a modest lift on the long-tail classes where
    the current bottleneck is Stage 3 retrieval, not Stage 1 detection.
  \item \textbf{Latency-budget arms.} A genuinely deployed system would
    need an A100 40GB / consumer GPU envelope, which constrains both
    Qwen-VL-7B and DINOv2-large. We have not measured the
    accuracy-vs-latency Pareto frontier under these constraints.
\end{itemize}

\paragraph{Why dish-correct (and not mAP) for the next pipeline.} We
believe the dish-level multiset framing generalises beyond cafeteria-tray
analysis to any product-grade detection task whose downstream consumer is
itemized rather than geometric: retail checkout, inventory-on-shelf
audits, kitchen prep verification. In all of these, an off-by-one count
error is a hard failure regardless of bbox quality. We hope that future
work in this lineage will report dish-correct (or its analogue) alongside
the customary IoU/mAP, so that the metric does not silently drift back
toward what is easy to measure rather than what matters.

\paragraph{Reproducibility.} The 80/10/10 split (seed $420$) is
deterministic and reproduced by
\texttt{stage1\_kcfd.dataset.\_compute\_stratified\_image\_split}.
The reference bank, prompt, evaluation set, and \texttt{score.py} metric
implementation are released alongside this paper.

\subsection{Limitations and threats to validity}
\label{sec:lim}

\paragraph{Sample size in the 3+\,item bucket.} Of the 596 held-out
images, only $7$ contain three or more items. The dish-correct rate on
that bucket ($28.6\%$ for the distilled arm, $0\%$ for the base) is
measured on a sample so small that any single error or correct prediction
shifts the rate by $\pm 14\,$pp. We report the number for completeness,
but our quantitative claims focus on the $1$- and $2$-item buckets where
$n \geq 100$. A future-work item is to deliberately curate a multi-item
evaluation supplement.

\paragraph{Reference bank leakage in spirit.} Both arms use the same
DINOv2 reference bank, built from train-split items only. We confirmed
that no held-out image's own crop appears as a reference. However, an
held-out image may share a \emph{kitchen, plate, lighting condition,
or camera angle} with its train neighbours, since v3.2 was collected at a
single cafeteria over a bounded period. Performance numbers in this
paper should therefore be read as \emph{within-distribution}; we make no
claim about generalisation to a different cafeteria, a different camera,
or a different time of year.

\paragraph{Class-prior bias in DINOv2 retrieval.} DINOv2 features are
trained on web-scale natural images. The frequency at which a given Saudi
food (e.g., \textit{masakaah}) appears in the DINOv2 pre-training data is
unknown but presumably very low; the retrieval is therefore biased toward
classes that are visually similar to common web images (rice, chicken)
and slightly disadvantaged on the visually-novel tail (asidah, hummus
cling-wrapped). A more food-specific encoder (e.g., DINOv3 fine-tuned on
the v3.2 reference set or a CLIP text-image model conditioned on the
class slug) is likely to recover several percentage points on the long
tail.

\paragraph{Single-cafeteria distribution.} v3.2 was collected at a
single university cafeteria. The 32 classes reflect that menu. We do
\emph{not} claim TahamajaNet will generalise zero-shot to another menu;
deploying at a new site would minimally require (i)~adding new classes
to the reference bank by photographing 5–10 reference images per item,
(ii)~optionally a brief fine-tune of Stage~1 on a few hundred new trays
to specialise the bbox style, and (iii)~refreshing the price table.

\paragraph{No human-rater check on dish-correct.} The ground-truth
multisets in the held-out set come from the v3.2 export, which itself
went through human cluster review at the \emph{cluster}, not per-image,
level. There is residual labeling noise (mis-clustered items, missed
items, ambiguous splits between rice and rice-with-sauce). A quick
ground-truth audit on a 30-image sample suggested $\sim$1--2\% of GT
multisets contain at least one disagreement with a fresh human rater.
We did not propagate this noise into our error bars.

\paragraph{Cardinality-aware retrain incomplete.} The structural count
loss curves in Figure~\ref{fig:training-curves} are from epoch~1 of the
in-progress 10-epoch retrain. The headline dish-correct number we report
($86.4\%$) is on the \texttt{epoch\_002} checkpoint of that retrain,
which is the strongest one available at submission. We expect a further
lift of several percentage points by epoch~$\geq 6$ but cannot report
it here.

\subsection{Broader impact and ethics}
\label{sec:ethics}

The pipeline exists to generate itemised cafeteria bills, replacing a
human cashier's keying-in step. The intended impact is reducing student
queue time at peak hours and freeing the cashier for higher-judgement
tasks (allergen checks, payment exceptions, complaint handling). We
explicitly do \emph{not} target unmanned operation: a cashier still
authorises every transaction, and the system surfaces its predicted
multiset as an editable list rather than auto-charging. We see two
genuine deployment risks: (i)~systematic mis-billing on rare classes
that costs students money (under-billing the cafeteria, over-billing the
student); and (ii)~the temptation to cut the human-in-the-loop step
once the queue-time gain is realised. We treat the human-in-the-loop
authorisation as a non-negotiable part of the deployment, not an
optional latency cost.

The v3.2 dataset contains photographs of food trays only. There are no
faces, no student IDs, no identifying information in the images. The
images were captured under an institutional agreement with the
cafeteria operator at KFUPM. We will release the dataset under a
research-only license that prohibits use for individual identification
or commercial sale.
